\documentclass{article}
\usepackage{amsmath}
\usepackage{geometry}
\geometry{a4paper, margin=1in}

\begin{document}

\section*{The Mathematics of 6DoF Optical Decoding}

To solve this, we don't try to calculate the physics of photons passing through slits. Instead, we use a \textbf{Linearized Calibration Matrix}. We assume that for very small movements, the change in light intensity is directly proportional to the movement of the top plate.

Let the 6DoF pose of the top plate relative to the bottom plate be represented by a 6x1 column vector $P$, containing translation $(x,y,z)$ and rotation $(\alpha, \beta, \gamma)$:
\begin{equation}
P = \begin{bmatrix} x \\ y \\ z \\ \alpha \\ \beta \\ \gamma \end{bmatrix}
\end{equation}

Let the raw brightness readings from your 6 photodetectors be represented by a 6x1 column vector $V$:
\begin{equation}
V = \begin{bmatrix} v_1 \\ v_2 \\ v_3 \\ v_4 \\ v_5 \\ v_6 \end{bmatrix}
\end{equation}

\subsection*{1. The Forward Model}
Assuming small displacements, the relationship between the physical pose $P$ and the sensor voltages $V$ is governed by a 6x6 sensitivity matrix (often called the Jacobian), denoted as $S$:
\begin{equation}
V = S \cdot P
\end{equation}

\subsection*{2. The Inverse Model (What your code actually does)}
To find the physical pose from your sensor readings, you multiply your voltage readings by the inverse of the sensitivity matrix. We call this the Calibration Matrix, $C$ (where $C = S^{-1}$):
\begin{equation}
P = C \cdot V
\end{equation}

Expanded, this looks like:
\begin{equation}
\begin{bmatrix} x \\ y \\ z \\ \alpha \\ \beta \\ \gamma \end{bmatrix} = \begin{bmatrix} 
c_{11} & c_{12} & c_{13} & c_{14} & c_{15} & c_{16} \\ 
c_{21} & c_{22} & c_{23} & c_{24} & c_{25} & c_{26} \\ 
c_{31} & c_{32} & c_{33} & c_{34} & c_{35} & c_{36} \\ 
c_{41} & c_{42} & c_{43} & c_{44} & c_{45} & c_{46} \\ 
c_{51} & c_{52} & c_{53} & c_{54} & c_{55} & c_{56} \\ 
c_{61} & c_{62} & c_{63} & c_{64} & c_{65} & c_{66} 
\end{bmatrix} \begin{bmatrix} v_1 \\ v_2 \\ v_3 \\ v_4 \\ v_5 \\ v_6 \end{bmatrix}
\end{equation}

\subsection*{3. Constructing the 4x4 Transformation Matrix ($T$)}
Once you have $P$, you map the translation $[x, y, z]^T$ directly into the $T$ matrix. You map the rotations $(\alpha, \beta, \gamma)$ by multiplying three standard rotation matrices together ($R = R_z(\gamma) R_y(\beta) R_x(\alpha)$):

\begin{equation}
T = \begin{bmatrix} 
R_{11} & R_{12} & R_{13} & x \\ 
R_{21} & R_{22} & R_{23} & y \\ 
R_{31} & R_{32} & R_{33} & z \\ 
0 & 0 & 0 & 1 
\end{bmatrix}
\end{equation}

\end{document}